% =====================================================================
%  Lecture 09 -- Evolved Packet Core (EPC)                (English edition)
% =====================================================================
\documentclass[11pt,a4paper]{article}
\providecommand\ACNroot{../..}
\input{\ACNroot/templates/lecture-en}

\ACNset{number}{09}
\ACNset{title}{Evolved Packet Core (EPC)}
\ACNset{shorttitle}{Evolved Packet Core}
\ACNset{titleother}{هستهٔ بستهٔ تکامل‌یافته (EPC)}
\ACNset{part}{Part II -- 4G access and core (Lectures 5--13)}
\ACNset{time}{7--9 hours (3.5 h reading, 4--5 h exercises)}
\ACNset{level}{Intermediate}
\ACNset{prereq}{Lectures 5--6; basic cryptography (symmetric keys, MAC); IP routing}
\ACNset{gen}{4}

\begin{document}
\ACNcover

% ---------------------------------------------------------------------
\section{Learning objectives}
After studying this lecture you should be able to:
\begin{objectives}
  \item \textbf{Explain} the design goals of the EPC and the division of work among MME, HSS,
        S-GW, P-GW and PCRF.
  \item \textbf{Describe} the EPS bearer model: default and dedicated bearers, TFTs and the QCI
        table, and \textbf{apply} GBR admission control.
  \item \textbf{Trace} EPS-AKA and derive the EPS key hierarchy from $K$ to the access-stratum keys.
  \item \textbf{Trace} PDN connectivity and dedicated bearer activation, modification and
        deactivation.
  \item \textbf{Describe} EMM/ECM states, tracking area updates and the EPC's role in handover.
  \item \textbf{Analyse} the impact of node failures and \textbf{estimate} the signalling storm after
        a gateway restart.
\end{objectives}

% ---------------------------------------------------------------------
\section{Prerequisite knowledge}
This lecture assumes the EPS architecture and reference points of Lecture~5 and the RRC/S1
procedures of Lecture~6. From security you need the idea of a shared symmetric key, of a message
authentication code (a keyed checksum that proves integrity and origin) and of a key derivation
function (KDF) that produces new keys from old ones without revealing them.

\paragraph{Key terms} The vocabulary of this lecture with its Persian equivalents:
\begin{keyterms}
PDN connection & \fa{اتصال PDN} & Association between a UE IP address and an APN, anchored in a P-GW. \\
Default bearer & \fa{حامل پیش‌فرض} & First bearer of each PDN connection; always on. \\
Dedicated bearer & \fa{حامل اختصاصی} & Additional bearer with its own QoS for selected flows. \\
TFT & \fa{قالب جریان ترافیک} & Packet filters mapping IP flows to bearers. \\
QCI & \fa{شناسهٔ کلاس کیفیت خدمات} & Index selecting priority, delay budget and loss rate. \\
EPS-AKA & \fa{توافق احراز هویت و کلید EPS} & Mutual authentication producing $K_{\mathrm{ASME}}$. \\
KDF & \fa{تابع استخراج کلید} & HMAC-SHA-256-based derivation of new keys. \\
EMM / ECM & \fa{مدیریت تحرک / اتصال EPS} & Registration state / signalling-connection state. \\
TAU & \fa{به‌روزرسانی ناحیهٔ ردیابی} & UE procedure when entering a TA outside its list or periodically. \\
PCC & \fa{کنترل سیاست و صورت‌حساب} & PCRF decisions enforced by the PCEF in the P-GW. \\
\end{keyterms}

% ---------------------------------------------------------------------
\section{Introduction and motivation}
The EPC was designed in the System Architecture Evolution (SAE) study with a clear brief: a
\emph{single}, \emph{all-IP} core for LTE that could also serve legacy 3GPP access (via SGSN) and
non-3GPP access (WLAN, CDMA2000), with lower latency, fewer nodes and a clean split between
control and user planes \cite{ts23401}. The resulting core has only five essential node types, yet
it must reproduce -- with packets -- everything the old CS and PS cores provided: subscriber
authentication, mobility with session continuity, QoS for voice, charging and lawful interception.

This lecture is the system view of the EPC. It connects the pieces: how a subscriber is
authenticated, how bearers are created and mapped to flows, how the UE's states are tracked, and
what happens when a node fails. Lectures~10--13 then zoom into the protocols and nodes (Diameter,
HSS, MME, gateways).

% ---------------------------------------------------------------------
\section{Fundamental concepts}

\subsection{Division of labour}
\begin{itemize}
  \item \textbf{MME} -- control plane only: NAS signalling, authentication, security, tracking areas,
        paging, bearer control, gateway selection.
  \item \textbf{HSS} -- subscription database and authentication centre.
  \item \textbf{S-GW} -- user-plane mobility anchor and idle-mode buffer.
  \item \textbf{P-GW} -- IP anchor towards each PDN, policy and charging enforcement (PCEF).
  \item \textbf{PCRF} -- policy and charging decisions (PCC rules).
\end{itemize}
The MME and the gateways are connected by S11 (GTPv2-C), the gateways by S5/S8; the MME talks to
the HSS over S6a and the P-GW to the PCRF over Gx, both Diameter.

\subsection{The bearer model}
\begin{definitionbox}{EPS bearer and TFT}
An EPS bearer is the unit of QoS in LTE: all packets on a bearer receive the same forwarding
treatment. A \emph{traffic flow template} (TFT) is a set of packet filters (source/destination
address and port ranges, protocol, type of service) that decides which bearer a packet uses: the
\emph{UL TFT} in the UE for uplink packets, the \emph{DL TFT} in the P-GW for downlink packets.
\end{definitionbox}
Each PDN connection starts with a default bearer without a TFT (it catches everything). Dedicated
bearers are created when a service needs different QoS -- typically triggered by the PCRF when an
application function (e.g.\ IMS) requests it. Every bearer is characterised by
\begin{itemize}
  \item \textbf{QCI} -- a standardised class with priority, packet delay budget (PDB) and packet error
        loss rate (PELR) (\cref{tab:qci});
  \item \textbf{ARP} -- allocation and retention priority, used for admission and pre-emption;
  \item for GBR bearers, \textbf{GBR} and \textbf{MBR} (guaranteed and maximum bit rate);
  \item for non-GBR bearers collectively, the \textbf{APN-AMBR} and \textbf{UE-AMBR} aggregate limits.
\end{itemize}

\begin{table}[htbp]
\centering\small
\caption{Standardised QCI characteristics (TS 23.203, Table 6.1.7-A, QCI 1--9) \cite{ts23203}.}\label{tab:qci}
\rowcolors{2}{acnmist}{white}
\begin{tabularx}{\linewidth}{@{}c c c c c Y@{}}
\toprule
\rowcolor{white}QCI & Type & Priority & PDB [ms] & PELR & Example services \\
\midrule
1 & GBR     & 2 & 100 & $10^{-2}$ & Conversational voice (VoLTE) \\
2 & GBR     & 4 & 150 & $10^{-3}$ & Conversational video (live) \\
3 & GBR     & 3 & 50  & $10^{-3}$ & Real-time gaming \\
4 & GBR     & 5 & 300 & $10^{-6}$ & Non-conversational (buffered) video \\
5 & non-GBR & 1 & 100 & $10^{-6}$ & IMS signalling \\
6 & non-GBR & 6 & 300 & $10^{-6}$ & Buffered video, TCP applications (premium) \\
7 & non-GBR & 7 & 100 & $10^{-3}$ & Voice, live video, interactive gaming \\
8 & non-GBR & 8 & 300 & $10^{-6}$ & TCP applications (e.g.\ web, e-mail) \\
9 & non-GBR & 9 & 300 & $10^{-6}$ & TCP applications, typical default bearer \\
\bottomrule
\end{tabularx}
\end{table}

\subsection{EMM and ECM states}
Two orthogonal state machines describe a UE in the EPC (\cref{fig:states}): \emph{EPS mobility
management} (EMM-DEREGISTERED/REGISTERED) says whether the network knows the UE; \emph{EPS
connection management} (ECM-IDLE/CONNECTED) says whether a NAS signalling connection (RRC + S1)
exists. A registered UE in ECM-IDLE keeps its IP address and bearers in the core (S5/S8 tunnels stay
up), but has no radio or S1-U resources; the S-GW buffers downlink data and asks the MME to page.

\begin{acnfigure}
  \acnfig{lecture-09/emm-ecm}
  \caption{EMM (registration) and ECM (signalling connection) states of a UE in the EPC.}
  \label{fig:states}
\end{acnfigure}

% ---------------------------------------------------------------------
\section{Detailed technical explanation}

\subsection{Authentication: EPS-AKA}
EPS-AKA \cite{ts33401} extends UMTS AKA (Lecture~3). The HSS generates an \emph{EPS authentication
vector} $(\text{RAND}, \text{XRES}, \text{AUTN}, K_{\mathrm{ASME}})$; the novelty is that
$K_{\mathrm{ASME}}$ is bound to the serving network identity, so a vector issued for one network is
useless in another. The MME sends RAND and AUTN; the USIM checks AUTN (which proves that the
challenge comes from the home network and is fresh, via the sequence number SQN), computes RES and
the same keys. The MME compares RES with XRES; both sides now share $K_{\mathrm{ASME}}$, which never
leaves the MME. The detailed Milenage computations are in Lecture~11.

\subsection{NAS and AS security}
Security is layered:
\begin{itemize}
  \item \textbf{NAS security} (UE--MME): ciphering and integrity protection of NAS messages with
        $K_{\mathrm{NASenc}}$ and $K_{\mathrm{NASint}}$, activated by the NAS \msg{Security Mode Command}.
  \item \textbf{AS security} (UE--eNB): integrity and ciphering of RRC ($K_{\mathrm{RRCint}}$,
        $K_{\mathrm{RRCenc}}$) and ciphering of user data ($K_{\mathrm{UPenc}}$) in PDCP, derived from
        $K_{\mathrm{eNB}}$ which the MME sends to the eNB in \msg{Initial Context Setup}.
\end{itemize}
User data is \emph{not} integrity-protected in LTE (only ciphered) -- a deliberate trade-off for
throughput that 5G revisits. The algorithms are EEA1/EIA1 (SNOW\,3G), EEA2/EIA2 (AES) and EEA3/EIA3
(ZUC); EEA0/EIA0 are null algorithms, EIA0 being allowed only for emergency calls without
authentication.

\subsection{Session management}
\paragraph{PDN connectivity} At attach (and later for additional APNs) the UE requests a PDN
connection. The MME selects a P-GW from the APN (DNS-based) and an S-GW from the UE's location, and
sends \msg{Create Session Request}; the P-GW allocates the IP address (IPv4, IPv6 prefix or both), asks
the PCRF for policy (Gx CCR-I/CCA-I) and creates the default bearer.
\paragraph{Dedicated bearer activation} Normally network-initiated (\cref{fig:dedicated}): the PCRF
provisions PCC rules, the P-GW sends \msg{Create Bearer Request} with TFT and QoS through the S-GW to
the MME, which sets up the E-RAB in the eNB and sends the NAS \msg{Activate Dedicated EPS Bearer
Context Request} to the UE.
\paragraph{Modification and deactivation} QoS or TFT changes use \msg{Update Bearer Request}; removal
uses \msg{Delete Bearer Request}. Deleting the default bearer deletes the whole PDN connection.

\subsection{Mobility in the EPC}
\begin{itemize}
  \item \textbf{Tracking area update (TAU)}: when the UE enters a TA not in its list, or when the
        periodic timer T3412 expires; the MME may change (context transfer from the old MME) and so may
        the S-GW.
  \item \textbf{X2 handover}: the core only switches the path (Lecture~4).
  \item \textbf{S1 handover}: used without X2 or when the MME or S-GW changes; the source and target
        eNBs exchange handover messages via the MME(s).
  \item \textbf{Inter-RAT mobility}: to UTRAN/GERAN through the SGSN (S3/S4 interfaces); the S-GW
        remains the anchor, the P-GW keeps the IP address.
\end{itemize}

\subsection{Failure scenarios}
\begin{itemize}
  \item \textbf{MME failure} -- with S1-flex the eNB connects to other MMEs of the pool; UEs whose
        contexts were lost are recovered when they next contact the network (TAU or service request is
        rejected with an implicit detach, followed by re-attach).
  \item \textbf{S-GW failure} -- detected through GTP echo and the \emph{restart counter} carried in the
        Recovery information element. The MME and P-GW delete the affected PDN connections; idle UEs are
        not reachable until they re-attach or the MME restores the sessions on another S-GW (restoration
        procedures).
  \item \textbf{P-GW failure} -- the IP anchor is lost: the UE's IP address and all its sessions on that
        APN disappear; applications must reconnect.
  \item \textbf{HSS unavailable} -- already registered UEs continue (the MME caches subscription and
        unused vectors); new attaches of UEs unknown to the MME fail.
\end{itemize}

% ---------------------------------------------------------------------
\section{Mathematical models}

\subsection{GBR admission control}\label{sec:admission}
\paragraph{Problem} How many GBR bearers can a cell or gateway admit?
\paragraph{Model} A resource (cell, link) has an average capacity $C$; the operator reserves a
fraction $\rho_{\mathrm{GBR}}$ for GBR traffic. Bearer $i$ requests $\mathrm{GBR}_i$. Admission
requires
\begin{equation}\label{eq:adm}
  \sum_{i\in\mathcal{A}}\mathrm{GBR}_i + \mathrm{GBR}_{\mathrm{new}} \le \rho_{\mathrm{GBR}}\,C ,
\end{equation}
otherwise the request is rejected, or a bearer with lower ARP is pre-empted if the new one may
pre-empt and the old one is pre-emptable.
\paragraph{Example} $C = \qty{30}{\mega\bit\per\second}$ average cell throughput, $\rho_{\mathrm{GBR}} = 0.5$,
VoLTE GBR $= \qty{40}{\kilo\bit\per\second}$: at most $15\times10^6/40\times10^3 = 375$ calls.
\paragraph{Limits} In a radio cell $C$ is not constant: it depends on where the users are (their MCS).
Radio admission therefore estimates the \emph{resource} (RBs) a bearer will need at its current
channel quality rather than its bit rate (\cref{exr:rbadm}).

\subsection{The EPS key hierarchy}\label{sec:keys}
\paragraph{Model} Keys are derived with the 3GPP KDF, HMAC-SHA-256 keyed with the parent key over a
string $S$ that encodes the purpose (an FC byte and parameters) \cite{ts33401}:
\begin{subequations}\label{eq:kdf}
\begin{align}
  K_{\mathrm{ASME}} &= \KDF\bigl(\mathrm{CK}\concat\mathrm{IK},\ \text{SN id},\ \mathrm{SQN}\xor\mathrm{AK}\bigr), \label{eq:kasme}\\
  K_{\mathrm{eNB}} &= \KDF\bigl(K_{\mathrm{ASME}},\ \text{uplink NAS COUNT}\bigr), \label{eq:kenb}\\
  K_{\mathrm{NASint}} &= \KDF\bigl(K_{\mathrm{ASME}},\ \text{alg.\ type},\ \text{alg.\ id}\bigr)\ \text{(truncated to 128 bits)}, \label{eq:knas}
\end{align}
\end{subequations}
and analogously for $K_{\mathrm{NASenc}}$, $K_{\mathrm{RRCint}}$, $K_{\mathrm{RRCenc}}$ and $K_{\mathrm{UPenc}}$
from $K_{\mathrm{eNB}}$ (\cref{fig:keys}).
\paragraph{Interpretation} Each key is bound to its context: the serving network (\cref{eq:kasme}),
the NAS COUNT, hence a fresh $K_{\mathrm{eNB}}$ at every transition to connected mode
(\cref{eq:kenb}), and the algorithm. Compromise of one eNB's keys reveals nothing about
$K_{\mathrm{ASME}}$ (one-way KDF). At handover a new $K_{\mathrm{eNB}}^\ast$ is derived so the target does
not learn the source's keys.
\paragraph{Brute force} A 128-bit algorithm key has $2^{128}\approx3.4\times10^{38}$ values; at $10^{12}$
trials per second an exhaustive search takes about $1.1\times10^{19}$ years. Security therefore rests on
protecting $K$ and on the algorithms, not on key length.

\subsection{Signalling storm after a gateway restart}\label{sec:storm}
\paragraph{Problem} An S-GW serving $N$ UEs restarts; all of them must re-establish their sessions.
How long does recovery take, and what load does it put on the MME?
\paragraph{Model} Each recovery costs $m$ messages at the MME; the MME processes at most $\mu$ messages
per second. If all UEs try immediately, the MME is saturated for
\begin{equation}\label{eq:storm}
  T_{\mathrm{rec}} = \frac{N\,m}{\mu}.
\end{equation}
If recovery attempts are spread uniformly over a window $T_\mathrm{w} > T_{\mathrm{rec}}$, the offered
signalling rate is $N m/T_\mathrm{w}$ and the MME stays below saturation.
\paragraph{Example} $N = 10^6$, $m = 10$, $\mu = 2\times10^4$ msg/s: $T_{\mathrm{rec}} = 500$\,s. Spreading
over $T_\mathrm{w} = 900$\,s gives $1.1\times10^4$ msg/s (56\% of capacity) instead of an instantaneous burst
of $10^7$ messages, which would overflow queues, time out, and trigger retries that make the storm worse.
\paragraph{Lesson} Restoration procedures and randomised timers exist to turn a burst into a spread
load -- the same positive-feedback hazard as RACH overload (Lecture~8), one level higher.

% ---------------------------------------------------------------------
\section{Network architecture and diagrams}
\Cref{fig:bearer} shows the bearer as seen end to end. The radio bearer, S1 bearer and S5/S8 bearer
are bound one-to-one; the S-GW and eNB simply relay packets between tunnels, while the TFTs at the two
ends decide the bearer. The E-RAB (radio + S1) is the part that is released in ECM-IDLE; the S5/S8
part persists.

\begin{acnfigure}
  \resizebox{\linewidth}{!}{\acnfig{lecture-09/bearer-structure}}
  \caption{EPS bearer structure with a default and a dedicated bearer. The UL TFT in the UE and the DL TFT
  in the P-GW map IP flows to bearers.}
  \label{fig:bearer}
\end{acnfigure}

\begin{acnfigure}
  \acnfig{lecture-09/key-hierarchy}
  \caption{EPS key hierarchy. Key lengths in bits; algorithm keys are truncated to 128 bits.}
  \label{fig:keys}
\end{acnfigure}

% ---------------------------------------------------------------------
\section{Step-by-step procedures}

\begin{procedurebox}{EPS-AKA and NAS security activation}\label{proc:aka}
\begin{steps}
  \item The UE sends \msg{Attach Request} with its GUTI (or IMSI if it has none).
  \item If no vector is available, the MME sends S6a \msg{Authentication Information Request} with the
        IMSI and the serving network identity; the HSS returns one or more EPS vectors.
  \item The MME sends \msg{Authentication Request} (RAND, AUTN, key set identifier).
  \item The USIM verifies AUTN (MAC and SQN freshness), computes RES, CK, IK; the UE derives
        $K_{\mathrm{ASME}}$ and returns \msg{Authentication Response} (RES).
  \item The MME checks RES = XRES and sends \msg{Security Mode Command} (selected algorithms, integrity
        protected with $K_{\mathrm{NASint}}$); the UE verifies it and answers \msg{Security Mode Complete},
        ciphered and integrity-protected.
\end{steps}
\end{procedurebox}

\begin{acnfigure}
  \resizebox{\linewidth}{!}{\acnfig{lecture-09/eps-aka}}
  \caption{EPS-AKA and NAS security mode (\cref{proc:aka}).}
  \label{fig:aka}
\end{acnfigure}

\begin{procedurebox}{Network-initiated dedicated bearer activation}\label{proc:ded}
\begin{steps}
  \item An application function (e.g.\ the IMS P-CSCF) requests resources over Rx; the PCRF decides
        PCC rules and pushes them to the P-GW over Gx.
  \item The P-GW binds the rule to a new bearer (QCI, ARP, GBR/MBR, DL TFT) and sends GTPv2-C
        \msg{Create Bearer Request} to the S-GW, which forwards it to the MME.
  \item The MME sends S1AP \msg{E-RAB Setup Request} carrying the NAS \msg{Activate Dedicated EPS Bearer
        Context Request} (UL TFT, QoS).
  \item The eNB admits the bearer, configures a radio bearer with \msg{RRC Connection Reconfiguration} and
        answers \msg{E-RAB Setup Response}; the UE answers the NAS request with \emph{Accept}.
  \item \msg{Create Bearer Response} travels back to the P-GW, which confirms to the PCRF (Gx RAA). Packets
        matching the TFTs now use the new bearer.
\end{steps}
\end{procedurebox}

\begin{acnfigure}
  \resizebox{\linewidth}{!}{\acnfig{lecture-09/dedicated-bearer}}
  \caption{Dedicated bearer activation triggered by policy (\cref{proc:ded}).}
  \label{fig:dedicated}
\end{acnfigure}

% ---------------------------------------------------------------------
\section{Worked examples}

\begin{examplebox}{Which bearer carries the packet?}\label{ex:tft}
\textbf{Problem.} A UE has a default bearer (QCI 9) and a dedicated bearer (QCI 1) whose DL TFT matches
UDP packets from 10.1.2.3, ports 40000--40100. Which bearer carries (a) an RTP packet from 10.1.2.3:40050,
(b) a TCP packet from 10.1.2.3:443, (c) a DNS response?

\textbf{Solution.} (a) Dedicated bearer (all filter components match). (b) Default bearer (protocol and
port do not match). (c) Default bearer (no filter matches; the default bearer has no TFT and receives
the rest).

\textbf{Interpretation.} Bearer selection is pure packet classification in the P-GW (downlink) and UE
(uplink); the radio then applies the QCI's treatment to the bearer.
\end{examplebox}

\begin{examplebox}{Admission by radio resources}\label{ex:rbadm}
\textbf{Problem.} A VoLTE bearer needs 40\,kbit/s. At its MCS a cell-centre UE carries 800 bits per RB pair
(one RB pair for one 1-ms subframe); a cell-edge UE only 160 bits. If 25 RB pairs per subframe are reserved for
GBR traffic, how many calls of each kind can be admitted?

\textbf{Solution.} RB pairs per subframe needed $= \mathrm{GBR}/(\text{bits per RB pair}\times1000\,\mathrm{s}^{-1})$:
centre $40\,000/(800\times1000) = 0.05$, edge $40\,000/(160\times1000) = 0.25$. Admissible: $25/0.05 =
500$ centre users or $25/0.25 = 100$ edge users.

\textbf{Interpretation.} The same GBR costs five times more radio resource at the cell edge -- admission
must be resource-based (\cref{sec:admission}).
\end{examplebox}

\begin{examplebox}{Key freshness}\label{ex:fresh}
\textbf{Problem.} A UE goes from idle to connected 200 times per day with the same $K_{\mathrm{ASME}}$.
How many different $K_{\mathrm{eNB}}$ values does it use, and why does that matter?

\textbf{Solution.} 200, because each transition uses a new uplink NAS COUNT in \cref{eq:kenb}. Even if the
same eNB serves all connections, PDCP COUNT values restart with each new key, so keystreams are never
reused under one key.
\end{examplebox}

% ---------------------------------------------------------------------
\section{Technology comparison}
\begin{table}[htbp]
\centering\small
\caption{Packet core functions: GPRS/UMTS PS core versus EPC versus 5GC.}\label{tab:core}
\rowcolors{2}{acnmist}{white}
\begin{tabularx}{\linewidth}{@{}l Y Y Y@{}}
\toprule
\rowcolor{white} & GPRS/UMTS PS & EPC & 5GC \\
\midrule
Mobility/session control & SGSN (both)            & MME (mobility), S-/P-GW-C (session) & AMF (mobility), SMF (session) \\
User-plane anchor        & GGSN                   & P-GW (S-GW as mobility anchor)      & UPF (one or more) \\
QoS unit                 & PDP context            & EPS bearer (QCI)                    & QoS flow (5QI) within a PDU session \\
IP at registration       & No (separate activation) & Yes (default bearer)             & Optional (PDU session separate) \\
Authentication           & UMTS AKA (quintets)    & EPS-AKA ($K_{\mathrm{ASME}}$)       & 5G-AKA / EAP-AKA$'$, home control \\
Subscriber database      & HLR (MAP)              & HSS (Diameter)                      & UDM/UDR (HTTP/2) \\
\bottomrule
\end{tabularx}
\end{table}

\paragraph{Trade-offs} The bearer model gives precise, per-bearer QoS but every QoS change needs a
signalling procedure through four nodes; 5G replaces it with QoS flows that can be added inside an
existing session with less signalling (Lecture~22). The MME's concentration of NAS, security and
bearer control is efficient but makes it a large stateful node; 5G splits it into AMF and SMF.

% ---------------------------------------------------------------------
\section{Key takeaways}
\begin{takeawaybox}
\begin{itemize}
  \item EPC: MME (control), HSS (subscribers), S-GW (mobility anchor), P-GW (IP anchor, PCEF), PCRF (policy).
  \item QoS unit = EPS bearer (QCI, ARP, GBR/MBR); TFTs map flows to bearers; default bearer always on.
  \item EPS-AKA binds $K_{\mathrm{ASME}}$ to the serving network; KDF chain $K\to$ CK,IK $\to K_{\mathrm{ASME}}\to K_{\mathrm{eNB}}\to$ AS keys.
  \item NAS security UE--MME, AS security UE--eNB; user data ciphered but not integrity-protected in LTE.
  \item EMM (registered?) and ECM (signalling connection?) are independent state machines.
  \item GBR admission must be resource-based; a gateway restart can cause a signalling storm of $Nm/\mu$ seconds
        unless recovery is spread.
\end{itemize}
\end{takeawaybox}

% ---------------------------------------------------------------------
\section{Exercises}

\begin{exercisebox}{1}{Idle UE}\label{exr:idle}
List which resources exist for a UE in EMM-REGISTERED + ECM-IDLE: IP address, S5/S8 tunnel, S1-U tunnel,
radio bearer, eNB context, MME context. What triggers their re-creation?
\end{exercisebox}

\begin{exercisebox}{1}{Serving-network binding}\label{exr:binding}
Explain what attack is prevented by including the serving network identity in the derivation of
$K_{\mathrm{ASME}}$.
\end{exercisebox}

\begin{exercisebox}{2}{Admission}\label{exr:rbadm}
Repeat \cref{ex:rbadm} for a mix of 70\% centre and 30\% edge users. How many VoLTE calls can be admitted?
\end{exercisebox}

\begin{exercisebox}{2}{Periodic TAU load}\label{exr:tau}
An MME serves $5\times10^6$ registered UEs with T3412 $=$ 54\,min. Each periodic TAU costs 6 messages.
Compute the average message rate. What happens to the load if T3412 is reduced to 12\,min?
\end{exercisebox}

\begin{exercisebox}{3}{Storm with retries}\label{exr:storm}
In \cref{sec:storm}, assume that requests waiting longer than 30\,s time out and are retried once. Model the
MME as a single queue with service rate $\mu$ and estimate the total number of messages processed when all
$N$ UEs try at once. Compare with spreading over $T_\mathrm{w} = 900$\,s.
\end{exercisebox}

\begin{exercisebox}{3}{Bearer versus flow}\label{exr:flows}
A UE runs a VoLTE call, a video call and background web traffic. Count the GTP tunnels on S1-U and S5 in LTE
(one PDN for IMS, one for the Internet) and discuss the signalling needed to add a second video stream. How
would a flow-based model reduce it?
\end{exercisebox}

\subsection*{Solutions}
\addcontentsline{toc}{subsection}{Solutions}

\begin{solutionbox}{\ref{exr:idle}}
Kept: IP address, S5/S8 tunnels, MME context (and S-GW/P-GW contexts). Released: S1-U tunnel, radio bearers,
eNB context. They are re-created by a Service Request -- triggered by the UE (uplink data or signalling) or by
paging after the S-GW receives downlink data and sends a Downlink Data Notification to the MME.
\end{solutionbox}

\begin{solutionbox}{\ref{exr:binding}}
Without binding, a vector (or key) obtained by one serving network could be used by another, e.g.\ a
compromised or rogue visited network could reuse vectors to impersonate a different network to the UE, or a
network could claim to be a cheaper one for billing. Because both the HSS and the UE derive
$K_{\mathrm{ASME}}$ with the SN id, a mismatch between the claimed network and the network that requested the
vector produces different keys and security mode fails.
\end{solutionbox}

\begin{solutionbox}{\ref{exr:rbadm}}
Mean cost per call $= 0.7\times0.05 + 0.3\times0.25 = 0.11$ RB pairs per ms; calls $= 25/0.11 \approx 227$.
\end{solutionbox}

\begin{solutionbox}{\ref{exr:tau}}
TAU rate $= 5\times10^6/(54\times60) = 1543$ per s; messages $= 6\times1543 = 9259$\,msg/s. With 12\,min the
rate increases by $54/12 = 4.5\times$ to about $41\,700$\,msg/s -- shorter timers detect lost UEs faster but
cost signalling and UE battery.
\end{solutionbox}

\begin{solutionbox}{\ref{exr:storm}}
Total work $Nm = 10^7$ messages at $\mu = 2\times10^4$\,msg/s needs 500\,s. Requests queued longer than 30\,s
time out: only the first $30\mu/m = 6\times10^4$ UEs complete in the first 30\,s; the others retry, adding up to
$N$ new attempts. A crude bound: with one retry each, offered work approaches $2Nm$ and recovery time about
1000\,s, with many UEs failing twice. Spreading over 900\,s keeps the queue short (load 56\%), no timeouts,
total work exactly $Nm$ -- faster overall despite the deliberate delay.
\end{solutionbox}

\begin{solutionbox}{\ref{exr:flows}}
IMS PDN: default bearer (QCI 5), dedicated QCI 1 (voice) and QCI 2 (video): 3 bearers. Internet PDN: default
bearer: 1. Total 4 EPS bearers, i.e.\ 4 S1-U tunnels and 4 S5 tunnels. A second video stream with the same QCI
can be added to the existing QCI 2 bearer by updating its TFT (Update Bearer procedure, still through PCRF,
P-GW, S-GW, MME, eNB and UE); a new QCI would need a new bearer. In 5G all flows of a PDU session share one
N3 tunnel; adding a QoS flow is a rule update in SMF, UPF and gNB with less per-flow tunnel state
(Lecture~22).
\end{solutionbox}

% ---------------------------------------------------------------------
\section{Further reading}
EPC procedures are specified in \cite{ts23401}, policy and charging control and the QCI table in
\cite{ts23203}, security in \cite{ts33401}, NAS in \cite{ts24301} and GTPv2-C in \cite{ts29274}. Olsson et
al.\ \cite{olsson2013epc} give a highly readable system description of the EPC.

\printbibliography[heading=subbibliography,title={\ACNstrRefs}]
\end{document}
